% =====================================================================
%  Capitulo 3 --- Planeamento
%  Traducao portuguesa de secoes/03-planeamento.tex
% =====================================================================
\chapter{Planeamento}

\section{Valores, missão e visão}
A AstraZeneca não usa os termos ``missão'' e ``visão''. Usa um \textit{Purpose} (propósito), que desempenha um papel semelhante ao da missão, ao explicar por que razão a organização existe, a \textit{Ambição 2030}, que corresponde à visão e fixa a direção pretendida, e cinco \textit{Valores}, que orientam a cultura organizacional e a tomada de decisão \cite[pp.~8, 10]{AR}, \cite[pp.~15--16]{PLAN}.

\subsection{Missão}
``Ultrapassamos as fronteiras da ciência para entregar medicamentos que mudam vidas'' \cite[p.~8]{AR}. A declaração é curta e clara quanto à razão de existir da empresa: criar medicamentos que mudem a vida dos doentes. O Grupo afirma ainda que pretende criar valor para os doentes, a sociedade, o planeta e os acionistas \cite[p.~8]{AR}, o que mostra que reconhece vários stakeholders e não apenas os seus acionistas. O âmbito não faz parte da missão, mas pode ser completado com o modelo de negócio: investigação, desenvolvimento, fabrico e comercialização de medicamentos sujeitos a receita médica em cuidados primários, cuidados especializados e doenças raras \cite[p.~8]{AR}.

\subsection{Visão}
Ser pioneira na ciência, liderar nas suas áreas de doença e transformar os resultados para os doentes, com metas concretas para 2030: pelo menos 20 medicamentos novos e \$80bn de Receita Total \cite[p.~10]{AR}. As metas numéricas tornam a visão mensurável e permitem acompanhar o progresso. Estas metas são usadas como objetivo estratégico na Secção~\ref{sec:objetivos}. A visão global é depois desdobrada em ambições para cada área terapêutica:

\begin{itemize}
  \item em Oncologia, ``eliminar o cancro como causa de morte'' \cite[p.~14]{AR};
  \item em BioPharmaceuticals, ``transformar os cuidados de milhares de milhões de pessoas que vivem com doenças crónicas'' \cite[p.~16]{AR};
  \item em Doenças Raras, ``abrir novas possibilidades'', com foco em doenças com necessidade médica significativa não satisfeita, onde o progresso científico foi ausente ou limitado \cite[p.~22]{AR}.
\end{itemize}

Cada unidade recebe, assim, a sua própria direção, alinhada com a visão do Grupo.

\subsection{Valores}

\begin{xltabular}{\linewidth}{@{}lY@{}}
\caption{Os valores da AstraZeneca. Fonte: \cite[p.~8]{AR}.}\label{tab:valores}\\*
\toprule
\textbf{Valor} & \textbf{Significado} \\*
\midrule
\endfirsthead
\toprule
\textbf{Valor} & \textbf{Significado} \\*
\midrule
\endhead
\bottomrule
\endfoot
Seguimos a ciência & Ultrapassar as fronteiras da ciência e trabalhar de forma criativa com parceiros \\
Colocamos os doentes primeiro & Compreender as necessidades dos doentes e tê-las em conta em cada decisão \\
Jogamos para ganhar & Construir equipas de alto desempenho, inclusivas e diversas, e fazer as escolhas certas para ganhar \\
Fazemos o que está certo & Aplicar padrões éticos elevados em todos os aspetos do negócio \\
Somos empreendedores & Agir com urgência, coragem e resiliência, assumindo riscos inteligentes \\
\end{xltabular}

O foco na ciência e nos doentes está diretamente ligado à investigação e ao desenvolvimento de medicamentos, enquanto a integridade e o espírito empreendedor sustentam uma conduta responsável e a procura de novas oportunidades. Os valores orientam a cultura e as decisões da empresa e são medidos no inquérito interno Pulse: 88\% dos colaboradores vêem a empresa como orientada para o doente \cite[p.~38]{AR}. Medi-los é uma forma de controlo, porque permite à empresa verificar se a cultura corresponde ao que se pretende.

\begin{analysis}
A missão, a visão e os valores são coerentes entre si e partilham a mesma ideia central: a ciência ao serviço do doente. O principal desafio para a gestão é conciliar o \textit{jogar para ganhar} com o \textit{colocar os doentes primeiro}, isto é, equilibrar a rentabilidade exigida pelos acionistas com o acesso dos doentes aos medicamentos, num contexto de pressão sobre os preços \cite[pp.~32--33]{AR}.
\end{analysis}

\section{Objetivos nos três níveis de gestão}\label{sec:objetivos}
Os objetivos estratégicos fixam os resultados globais de longo prazo, os objetivos táticos traduzem a estratégia em planos departamentais e os objetivos operacionais orientam atividades concretas do dia a dia \cite[pp.~5--9]{PLAN}. A AstraZeneca divulga apenas os seus objetivos estratégicos. Os objetivos táticos e operacionais abaixo são aqui propostos, não divulgados pela empresa, desenhados de forma a que cada um contribua para o nível acima. A Ambição 2030 é perseguida através da estratégia \textit{Growth Through Innovation} (crescimento através da inovação), que tem três prioridades: Ciência e Inovação, Crescimento e Liderança em Áreas Terapêuticas, e Pessoas e Sustentabilidade \cite[p.~10]{AR}. O objetivo tático abaixo insere-se na segunda prioridade e o objetivo operacional apoia a primeira.

\begin{xltabular}{\linewidth}{@{}N{3.1cm}Yl@{}}
\caption{Um objetivo para cada nível de gestão. Só o objetivo estratégico é declarado pelo Grupo.}\label{tab:objetivos}\\*
\toprule
\textbf{Nível} & \textbf{Objetivo} & \textbf{Horizonte} \\*
\midrule
\endfirsthead
\toprule
\textbf{Nível} & \textbf{Objetivo} & \textbf{Horizonte} \\*
\midrule
\endhead
\bottomrule
\endfoot
\textbf{Estratégico}\newline CEO e SET & Lançar pelo menos 20 medicamentos novos e alcançar \$80bn de Receita Total até 2030 \cite[p.~10]{AR} & Longo \\
\textbf{Tático}\newline Diretor de desenvolvimento clínico em oncologia & Submeter pelo menos três novos medicamentos de oncologia a aprovação regulamentar até ao final de 2027 & Médio \\
\textbf{Operacional}\newline Chefe de equipa de um ensaio clínico & Revisão de pelo menos 98\% dos documentos do ensaio dentro do prazo, sem achados críticos, até ao final do próximo trimestre & Curto \\
\end{xltabular}

\textbf{Estratégico.} Abrange toda a empresa, tem um horizonte longo e exige decisões que só a gestão de topo pode tomar, como o investimento em I\&D, que áreas priorizar e as aquisições.

\textbf{Tático.} Transforma a meta estratégica num resultado para uma unidade específica, a Oncologia, que é a maior área (44\% da receita) \cite[p.~2]{AR} e teve 10 resultados positivos de Fase III em 2025 \cite[p.~13]{AR}. Pertence a um gestor intermédio, que coordena equipas mas não define a estratégia do Grupo.

\textbf{Operacional.} Centra-se numa tarefa concreta, tem um horizonte curto e um indicador preciso. A qualidade da documentação é essencial para a aprovação regulamentar, e o Grupo já reporta zero observações críticas em 42 inspeções externas \cite[p.~4]{AR}.

\begin{analysis}
Os três objetivos são SMART e formam uma cascata. Documentação fiável (operacional) permite submissões regulamentares (tático), e as submissões aprovadas tornam-se medicamentos novos e receita (estratégico).
\end{analysis}

\section{Análise SWOT e estratégia}
A SWOT é construída olhando primeiro para fora da empresa e depois para dentro, seguindo o material da unidade curricular \cite[pp.~17--29]{PLAN}. Fora, a PESTEL cobre a envolvente geral \cite[pp.~19--21]{PLAN} e as cinco forças de Porter cobrem o setor farmacêutico \cite[pp.~21--24]{PLAN}. Dentro, a cadeia de valor mostra o que a empresa faz \cite[p.~25]{PLAN} e a VRIO mostra que recursos lhe dão vantagem \cite[p.~26]{PLAN}. As forças e as fraquezas descrevem a empresa hoje; as oportunidades e as ameaças descrevem o que está a chegar \cite[p.~27]{PLAN}. Os factos vêm do Relatório e Contas e de outras fontes citadas, enquanto as classificações e os graus atribuídos são uma avaliação, não declarações da empresa.

\begin{figure}[H]\centering
\begin{tikzpicture}[node distance=7mm and 6mm]
\node[fl] (a1) {PESTEL\\{\color{azmuted}envolvente geral}};
\node[fl, right=of a1] (a2) {Cinco forças\\{\color{azmuted}o setor}};
\node[fl, right=of a2] (a3) {Oportunidades \& Ameaças\\{\color{azmuted}fora, futuro}};
\node[fl, below=of a1] (b1) {Cadeia de valor\\{\color{azmuted}atividades}};
\node[fl, right=of b1] (b2) {VRIO\\{\color{azmuted}recursos}};
\node[fl, right=of b2] (b3) {Forças \& Fraquezas\\{\color{azmuted}dentro, hoje}};
\node[fle, right=of a3] (sw) {SWOT};
\node[fle, right=of b3] (st) {Cruzamento\\$\rightarrow$ Estratégia};
\draw[se] (a1)--(a2); \draw[se] (a2)--(a3); \draw[se] (a3)--(sw);
\draw[se] (b1)--(b2); \draw[se] (b2)--(b3); \draw[se] (b3)--(st);
\draw[se] (sw)--(st);
\end{tikzpicture}
\caption{Da envolvente à estratégia, seguindo o método do material da unidade curricular \cite[pp.~17--29]{PLAN}.}
\label{fig:metodo}
\end{figure}
\subsection{A envolvente geral: PESTEL}
O cenário fora da AstraZeneca é misto. A procura está do lado da empresa: as pessoas vivem mais tempo, as doenças crónicas e o cancro continuam a crescer, e a despesa em saúde continua a subir. O preço e a regulação trabalham contra ela: os orçamentos de saúde estão apertados, os EUA passaram a negociar preços, e os governos querem limitar a proteção das patentes. O material da unidade curricular pede que se classifique o impacto de cada fator e que se olhe para a forma como os fatores interagem \cite[p.~19]{PLAN}; a Tabela~\ref{tab:pestel} apresenta isso fator a fator.


\begin{xltabular}{\linewidth}{@{}N{2.6cm}Yp{2.45cm}p{1.5cm}@{}}
\caption{Análise PESTEL. A inflação, os níveis de escolaridade e a estabilidade governamental ficam de fora porque o Relatório e Contas não os trata como fatores de negócio.}\label{tab:pestel}\\*
\toprule
\textbf{Fator} & \textbf{Evidência} & \textbf{Classe} & \textbf{Impacto} \\*
\midrule
\endfirsthead
\toprule
\textbf{Fator} & \textbf{Evidência} & \textbf{Classe} & \textbf{Impacto} \\*
\midrule
\endhead
\bottomrule
\endfoot
\multicolumn{4}{@{}l@{}}{\textbf{Político}} \\*
\midrule
Geopolítica & Tensão e rivalidade estratégica entre os EUA e a China, e nacionalismo económico crescente \cite[p.~6]{AR}. Controlos ao comércio, tarifas e impostos podem aumentar custos ou reduzir receita \cite[p.~49]{AR}. & \bbad{Ameaça} & Alto \\
Localização e tarifas & Os governos incentivam as empresas a localizar operações para reduzir o risco da cadeia de fornecimento \cite[p.~7]{AR}. O Grupo obteve uma isenção de três anos às tarifas da Secção 232 e planeia investir \$50bn em fabrico e I\&D nos EUA \cite[pp.~4, 32--33]{AR}. & \bgood{Oportunidade} & Médio \\
Fiscalidade & O Grupo pagou \$2,845m de imposto sobre o rendimento em 2025. A taxa de imposto Core foi de 18\% e está prevista entre 18--22\% em 2026 \cite[pp.~59, 64]{AR}. & \bneut{Neutro} & Baixo \\
\midrule
\multicolumn{4}{@{}l@{}}{\textbf{Económico}} \\*
\midrule
Orçamentos e preços & Os orçamentos da saúde estão sob pressão, o que leva a pressão descendente sobre os preços \cite[p.~7]{AR}. A consolidação dos pagadores e os reembolsos restritivos geram pressão significativa sobre os preços \cite[p.~32]{AR}. & \bbad{Ameaça} & Alto \\
Crescimento do mercado & Espera-se que a despesa global em saúde cresça 7.1\% por ano em 2025--2029, e as vendas farmacêuticas globais cresceram 10.0\% em 2025 \cite[p.~6]{AR}. A previsão de crescimento do PIB global é de 3.3\% em 2026 e 3.2\% em 2027 \cite[p.~7]{AR}. & \bgood{Oportunidade} & Alto \\
Deslocação para mercados emergentes & O poder económico está a passar do G7 para grandes mercados emergentes como a China, a Índia e a Indonésia \cite[p.~6]{AR}. & \bgood{Oportunidade} & Médio \\
Taxas de juro e capital & Os empréstimos ao abrigo de linhas contratadas são indexados à SOFR mais uma margem, e a tesouraria gere o risco de taxa de juro, de câmbio e de crédito \cite[pp.~61, 64]{AR}. O Grupo dispõe de \$4,875m de linhas contratadas não utilizadas e sem covenants \cite[p.~61]{AR}. & \bneut{Neutro} & Baixo \\
Câmbios & Alterações significativas nas taxas de câmbio podem afetar negativamente o negócio \cite[p.~49]{AR}. Às taxas de janeiro de 2026, a Receita Total de 2026 ganharia uma percentagem baixa de um dígito \cite[p.~64]{AR}. & \bneut{Neutro} & Baixo \\
\midrule
\multicolumn{4}{@{}l@{}}{\textbf{Social}} \\*
\midrule
Envelhecimento & A população com 65 ou mais anos deverá triplicar até 2100, o que aumenta a procura de cuidados \cite[pp.~6--7]{AR}. & \bgood{Oportunidade} & Alto \\
Carga de doença & As doenças não transmissíveis representam 75\% das mortes não pandémicas. Cerca de 1.4 mil milhões de pessoas têm hipertensão, estimam-se mais de 30 milhões de novos diagnósticos de cancro até 2040, e menos de 10\% das doenças raras têm um tratamento aprovado \cite[pp.~7, 13, 17, 23]{AR}. & \bgood{Oportunidade} & Alto \\
Talento & O Grupo desenvolve talento internamente: 5,100 promoções e 2,760 contratações externas nos seus centros globais em 2025 \cite[p.~39]{AR}. & \bneut{Neutro} & Baixo \\
\midrule
\multicolumn{4}{@{}l@{}}{\textbf{Tecnológico}} \\*
\midrule
Inteligência artificial & Estima-se que a IA generativa crie \$4--7bn de valor por ano para a indústria farmacêutica, e o AR nota que a IA ``vai libertar valor e trazer novos riscos'' \cite[p.~7]{AR}. & Oportunidade com risco & Médio \\
Capacidade de I\&D e colaboração & A despesa em I\&D foi de \$14,232m e estão em curso mais de 100 estudos de Fase III \cite[pp.~3, 55]{AR}. Mais de 1,000 colaborações dão acesso a tecnologia \cite[p.~37]{AR}. & \bgood{Oportunidade} & Médio \\
\midrule
\multicolumn{4}{@{}l@{}}{\textbf{Ambiental}} \\*
\midrule
Clima & A crise climática agrava as desigualdades em saúde e pressiona os sistemas de saúde \cite[p.~7]{AR}. O Grupo tem metas próprias com base científica, incluindo zero emissões líquidas até 2045 \cite[p.~42]{AR}. Os fenómenos climáticos extremos ameaçam instalações, fornecedores e rotas de distribuição \cite[p.~44]{AR}. & \bbad{Ameaça} & Médio \\
Água, resíduos e natureza & As operações reduziram o consumo de água em 24\% e os resíduos em 23\% face a 2015 \cite[p.~33]{AR}. O Grupo reconhece os princípios ativos dos seus produtos como o seu impacto mais material na poluição da água \cite[p.~45]{AR}. & \bbad{Ameaça} & Baixo \\
\midrule
\multicolumn{4}{@{}l@{}}{\textbf{Legal}} \\*
\midrule
Regulação de preços (IRA) & Nos EUA, o Inflation Reduction Act sujeita medicamentos a preços negociados pelo Medicare. A Farxiga foi selecionada e a Calquence entra em vigor em 2027. O AR espera que o impacto na Farxiga seja gerível \cite[p.~33]{AR}; aqui é classificado como alto, porque a Farxiga é o maior produto, com \$8,492m, 14.5\% da Receita Total \cite[p.~56]{AR}. & \bbad{Ameaça} & Alto \\
Propriedade intelectual & Governos e pagadores pressionam para limitar a proteção da propriedade intelectual \cite[p.~49]{AR}. & \bbad{Ameaça} & Alto \\
Reforma farmacêutica da UE & O acordo provisório de dezembro de 2025 fixa oito anos de proteção de dados mais um ano de proteção de mercado (prorrogável por um ano se duas de três condições se verificarem), e uma exceção Bolar para que os genéricos possam estar disponíveis no primeiro dia após a expiração da propriedade intelectual \cite{EU}. & \bbad{Ameaça} & Médio \\
Regulação do setor & Atrasos nas aprovações regulamentares podem atrasar a comercialização e reduzir a receita \cite[p.~48]{AR}. Os produtos precisam de aprovação de autoridades como a FDA e a EMA \cite[p.~20]{AR}. & \bbad{Ameaça} & Médio \\
Regras de dados e IA & Um conjunto crescente de exigências legislativas e regulamentares abrange os dados e a IA \cite[p.~48]{AR}. & \bbad{Ameaça} & Baixo \\
\end{xltabular}

\begin{analysis}
Os fatores não agem isoladamente. As tarifas dos EUA, o acordo de preços e o fabrico local vêm como um pacote único: uma isenção tarifária em troca de preços alinhados com os de outros países ricos e de medicamentos fabricados nos EUA \cite[pp.~4, 32--33]{AR}. Em 2026 o preço IRA da Farxiga chega no mesmo ano em que esta perde a exclusividade nos EUA, embora o AR espere que o impacto seja gerível \cite[p.~33]{AR}. O envelhecimento empurra a procura para cima enquanto os orçamentos apertados seguram os preços em baixo, e é por isso que o Grupo usa preços diferenciados por patamares \cite[pp.~6--7, 41]{AR}. Olhando para a frente, a despesa em saúde continua a crescer até 2029, a Calquence entra em negociação de preço em 2027 e a meta de zero emissões líquidas estende-se até 2045 \cite[pp.~6, 33, 42]{AR}.
\end{analysis}

\subsection{O setor: as cinco forças de Porter}
Esta secção olha para o mercado dos medicamentos inovadores sujeitos a receita médica. Cada força é classificada como forte, moderada ou fraca \cite[p.~21]{PLAN}, usando os fatores que o material da unidade curricular enumera para cada uma \cite[pp.~22--24]{PLAN}.


\begin{xltabular}{\linewidth}{@{}N{2.5cm}Yp{2.3cm}@{}}
\caption{As cinco forças de Porter. A intensidade não é classificada no Relatório e Contas.}\label{tab:porter}\\*
\toprule
\textbf{Força} & \textbf{Evidência} & \textbf{Intensidade} \\*
\midrule
\endfirsthead
\toprule
\textbf{Força} & \textbf{Evidência} & \textbf{Intensidade} \\*
\midrule
\endhead
\bottomrule
\endfoot
\textbf{Rivalidade} & \textit{Número e equilíbrio:} vários grupos grandes de dimensão semelhante. Em 2024 os 10 maiores obtiveram entre \$44.47bn (Sanofi) e \$64.17bn (Merck \& Co.), e a AstraZeneca foi 5.ª com \$54.07bn, atrás da Merck \& Co., Pfizer, J\&J e AbbVie \cite{P50}. \textit{Custos fixos:} I\&D de \$14,232m por ano \cite[p.~55]{AR} e uma rede de fabrico de 31 instalações \cite[p.~33]{AR}. \textit{Diferenciação:} produtos protegidos por patentes e posições de liderança até ao fim da exclusividade \cite[pp.~9, 56]{AR}. \textit{Crescimento:} o setor está a crescer, com a despesa a subir 7.1\% por ano até 2029 e as vendas a subir 10.0\% em 2025 \cite[p.~6]{AR}, o que alivia a rivalidade. Estão em curso mais de 100 estudos de Fase III \cite[p.~3]{AR}. & \bbad{Forte} \\
\textbf{Novos concorrentes} & \textit{Capital e tempo:} estima-se que desenvolver um medicamento custe \$2.6bn (em dólares de 2013) \cite{MULLARD}, e a I\&D leva 5--15 anos e precisa de aprovação regulamentar \cite[p.~9]{AR}. \textit{Escala:} vendas em mais de 125 países com 47,400 colaboradores na área comercial \cite[p.~32]{AR}, e mais de 1,000 colaborações \cite[p.~37]{AR}. \textit{Biotecnológicas:} o Grupo reforça o pipeline através de aquisições, licenciamento e colaborações \cite[p.~48]{AR}, por exemplo a EsoBiotec \cite[p.~37]{AR}. \textit{Condição:} o AR avisa que limites à proteção da propriedade intelectual levariam a ``concorrência acelerada e intensificada'' \cite[p.~49]{AR}. \textit{Barreiras à saída:} o Grupo opera instalações de fabrico especializadas, e o AR regista o fecho de uma fábrica de comprimidos (Bangalore, maio de 2025) \cite[p.~33]{AR}. & \bgood{Fraca}\newline{\scriptsize condicional à PI} \\
\textbf{Fornecedores} & \textit{Concentração:} a maior parte do fabrico de API é externa, complementada por capacidade interna crescente \cite[p.~33]{AR}. A Índia representa 48\% dos Drug Master Files de API ativos submetidos à FDA dos EUA e a China 18\% em 2024; os ficheiros refletem a localização do fabricante, não o volume \cite{USP}. \textit{Substitutos:} alguns inputs não têm alternativa fácil, como os reagentes de sangue de caranguejo-ferradura até à adoção de uma alternativa sintética em 2025 \cite[p.~45]{AR}. \textit{Mitigação:} dupla fonte de fornecimento, múltiplos fornecedores e monitorização de stocks \cite[p.~48]{AR}. & \bwarn{Moderada} \\
\textbf{Clientes} & \textit{Concentração:} os pagadores são sobretudo seguradoras e orçamentos nacionais. A consolidação dos pagadores e os reembolsos restritivos geram pressão significativa sobre os preços, e aplicam-se descontos \cite[p.~32]{AR}. Nos EUA, três pharmacy benefit managers (CVS Caremark, Express Scripts e Optum Rx) processaram cerca de 80\% dos pedidos de receita equivalentes em 2025 \cite{FEIN}. Os preços são negociados ao abrigo do IRA \cite[p.~33]{AR} e seguem uma política de quatro patamares baseada no rendimento nacional \cite[p.~41]{AR}. Os produtos são prescritos por médicos, que não os pagam \cite[p.~32]{AR}. A pressão sobre os preços é um Risco Principal \cite[p.~48]{AR}. & \bbad{Forte} \\
\textbf{Produtos substitutos} & \textit{Custo:} quando a exclusividade termina, os genéricos entram a preços mais baixos. As vendas da Brilinta caíram 38\% para \$823m após a entrada de genéricos nos EUA e na Europa \cite[pp.~18, 56]{AR}, e a Symbicort erodiu na UE e no Japão \cite[p.~20]{AR}. A Farxiga ainda cresceu 10\% em 2025 apesar da concorrência de genéricos em alguns mercados \cite[p.~56]{AR}. A IQVIA estima mais de \$90bn em vendas que vão perder exclusividade em 2025--2029 \cite{IQVIA}. Os genéricos são aqui tratados como produtos substitutos, porque substituem o produto de marca pelo mesmo princípio ativo a um preço mais baixo \cite[p.~24]{PLAN}. & \bbad{Forte}\newline{\scriptsize ao longo do tempo, por produto} \\
\end{xltabular}

\begin{analysis}
Para quem já está dentro, o setor é atrativo: é muito difícil entrar. Mas cada produto só rende bem durante um tempo limitado. Enquanto um medicamento está protegido, os pagadores têm a vantagem; depois de a patente expirar, os genéricos tomam conta. Por isso, o verdadeiro jogo é renovar continuamente o portefólio. Duas classificações dependem de onde se traça a fronteira: se os genéricos fossem contados como novos concorrentes, essa força seria mais forte, e a baixa ameaça de entrada só se mantém enquanto as patentes continuarem protegidas \cite[p.~49]{AR}.
\end{analysis}

\subsection{Dentro da empresa: a cadeia de valor}
A cadeia de valor divide o que a AstraZeneca faz em atividades principais, que produzem e entregam os medicamentos, e atividades de suporte, que tornam o resto possível. O material da unidade curricular liga ainda estas atividades à estrutura de custos e ao lucro \cite[p.~25]{PLAN}, e pergunta o que a empresa faz bem e menos bem, incluindo as suas pessoas e a sua cultura \cite[p.~26]{PLAN}.


\begin{table}[H]\centering
\caption{Para onde vai cada dólar de receita em 2025, em \% da Receita Total, calculado a partir da demonstração de resultados \cite[p.~125]{AR}. A margem bruta é de 81.9\%.}
\label{tab:custos}
\begin{tabular}{@{}lrB@{}}
\toprule
\textbf{Rubrica} & \textbf{\% da receita} & \\
\midrule
Comerciais e administrativos & 33.9 & \barra{1.000} \\
I\&D                        & 24.2 & \barra{0.714} \\
Custo das vendas            & 18.1 & \barra{0.534} \\
Distribuição                &  1.0 & \barra{0.029} \\
\midrule
\textbf{Margem operacional}  & \textbf{23.4} & \barrab{0.690} \\
\bottomrule
\end{tabular}
\end{table}
\needspace{16\baselineskip}
\begin{xltabular}{\linewidth}{@{}N{2.45cm}YYp{2.45cm}@{}}
\caption{Análise da cadeia de valor. As classificações não são dadas no Relatório e Contas.}\label{tab:cadeia-valor}\\*
\toprule
\textbf{Atividade} & \textbf{Evidência} & \textbf{Pontos fracos} & \textbf{Avaliação} \\*
\midrule
\endfirsthead
\toprule
\textbf{Atividade} & \textbf{Evidência} & \textbf{Pontos fracos} & \textbf{Avaliação} \\*
\midrule
\endhead
\bottomrule
\endfoot
\multicolumn{4}{@{}l@{}}{\textbf{Atividades principais}} \\*
\midrule
Logística de entrada & A maior parte do fabrico de API é externa, complementada por capacidade interna crescente, e um programa de financiamento da cadeia de fornecimento apoia os fornecedores externos \cite[p.~33]{AR}. & A falta de medicamentos atingiu níveis recorde em 2025 \cite[p.~33]{AR}; não conseguir manter o fornecimento é um Risco Principal \cite[p.~48]{AR}. & \bwarn{Neutro / risco} \\
Operações & 31 instalações de operações em 15 países (excluindo as de fornecimento clínico) e desempenho de fornecimento superior a 99\% \cite[pp.~4, 33]{AR}. 63 inspeções GMP e GDP de autoridades de saúde com zero achados críticos \cite[p.~31]{AR}. A revisão do CEO cita zero recolhas ao nível do doente e zero observações críticas em 42 inspeções externas \cite[p.~4]{AR}. & 10 recolhas de produto \cite[p.~31]{AR}. O AR não reconcilia as contagens de 42 e de 63 inspeções. & \bgood{Força} \\
Logística de saída & A despesa de distribuição foi de \$579m, 1.0\% da receita (calculado) \cite[p.~125]{AR}. São reportados 217 lançamentos no mercado dentro do prazo em Operações \cite[pp.~27, 33]{AR}. & O AR não dá dados de desempenho de entrega. & \bneut{Neutro} \\
Marketing e vendas & 47,400 colaboradores na área comercial \cite[p.~27]{AR}; 16 medicamentos acima de \$1bn \cite[p.~4]{AR}; vendas através de empresas de marketing detidas integralmente, distribuidores e escritórios locais \cite[p.~32]{AR}. & Cinco produtos são 51\% da receita (calculado: \$30,045m ÷ \$58,739m) e a Farxiga sozinha é 14.5\% \cite[p.~56]{AR}. Os EUA são 43\% da receita, com pressão sobre os preços \cite[pp.~2, 32]{AR}. & \bwarn{Força} \newline{\scriptsize risco de concentração} \\
Serviços pós-venda & Farmacovigilância ao longo de todo o ciclo de vida do produto \cite[p.~31]{AR}; programas de apoio ao doente \cite[p.~41]{AR}; mais de 200 parcerias com sistemas de saúde em 50 países \cite[p.~4]{AR}. & 10 recolhas de produto \cite[p.~31]{AR}. & \bgood{Força} \\
\midrule
\multicolumn{4}{@{}l@{}}{\textbf{Atividades de suporte}} \\*
\midrule
I\&D e tecnologia & Despesa de I\&D de \$14,232m, 24.2\% da receita (calculado) \cite[pp.~28, 125]{AR}; 197 projetos em pipeline \cite[p.~30]{AR}; mais de 90\% da descoberta de moléculas pequenas com apoio de IA \cite[p.~36]{AR}. & 20 projetos descontinuados em 2025, 11 por segurança ou eficácia \cite[p.~30]{AR}; o ensaio de Fase III da Fasenra na DPOC não atingiu o objetivo primário \cite[p.~20]{AR}. & \bgood{Força} \\
Recursos humanos e cultura & 86\% dizem que a AstraZeneca é um ótimo lugar para trabalhar (inquérito Pulse) \cite[pp.~5, 38]{AR}; 89\% reportam formação ou desenvolvimento \cite[p.~39]{AR}; mais de 50,000 colaboradores participaram no programa de IA \cite[pp.~5, 39]{AR}. Os cinco valores do Grupo estão em \cite[p.~8]{AR}. & 4,441 incumprimentos do Código de Ética registados e 434 pessoas afastadas das suas funções \cite[p.~34]{AR}. A rotatividade não é divulgada. & \bgood{Força} \\
Gestão financeira & Fluxo de caixa operacional de \$14,575m (+23\%) \cite[p.~10]{AR}. Prioridades de capital: investir no negócio e no pipeline, manter um rating de grau de investimento, prosseguir o desenvolvimento de negócio e pagar um dividendo progressivo \cite[p.~10]{AR}. & Dívida líquida de \$23,374m, 1.6 vezes o fluxo de caixa operacional (calculado) \cite[p.~61]{AR}. & \bwarn{Força} \newline{\scriptsize alavancagem a monitorizar} \\
Recursos físicos & Uma fábrica de \$4.5bn na Virgínia e uma expansão de \$2bn em Maryland \cite[pp.~4, 33]{AR}. As emissões de Âmbito 1 e 2 estão 88.1\% abaixo de 2015 \cite[p.~42]{AR}. & Face a uma meta de $-$98\% até 2026 \cite[p.~42]{AR}. A fábrica de comprimidos de Bangalore fechou em maio de 2025 \cite[p.~33]{AR}. & \bgood{Força} \\
Aprovisionamento & Dupla fonte de fornecimento, múltiplos fornecedores e monitorização de stocks \cite[p.~48]{AR}. & O AR descreve a mitigação, não o desempenho. & \bneut{Neutro} \\
Administração e risco & Cada função do SET avalia o risco trimestralmente, alimentando o Group Risk Report para o Conselho de Administração, a Comissão de Auditoria e o SET \cite[p.~47]{AR}. O programa de reestruturação PAAGR gerou benefícios anuais recorrentes antes de impostos de \$1.9bn \cite[p.~59]{AR}. & Todas as decisões operacionais significativas são tomadas ao nível do SET \cite[p.~147]{AR} & \bwarn{Neutro / risco} \\
\midrule
\multicolumn{4}{@{}l@{}}{\textbf{Margem}} \\*
\midrule
Estrutura de custos e margem & 2025, \% da Receita Total (calculado a partir de \cite[p.~125]{AR}): custo das vendas 18.1\%, distribuição 1.0\%, I\&D 24.2\%, custos comerciais, gerais e administrativos 33.9\%. Margem bruta 81.9\% e margem operacional 23.4\%. & --- & \bgood{Força} \\
\end{xltabular}

\begin{analysis}
A maior parte da cadeia funciona bem: o fornecimento é fiável, a despesa em investigação é elevada e o negócio gera bastante liquidez. Os pontos fracos estão nas extremidades: as matérias-primas dependem de fornecedores externos, as vendas dependem de um punhado de produtos, e alguns projetos de investigação falham. O material da unidade curricular pede que se compare com os concorrentes \cite[p.~24]{PLAN}, mas a única comparação com pares no Relatório e Contas é o retorno para o acionista a dez anos: 307\%, contra 197\% dos pares farmacêuticos globais e 65\% dos pares europeus \cite[p.~90]{AR}. Para tudo o resto, o Grupo só pode ser comparado com o seu próprio histórico e com as suas metas.
\end{analysis}

\subsection{Que recursos dão vantagem: VRIO}
A VRIO faz quatro perguntas sobre cada recurso: cria valor, é raro no setor, é difícil ou dispendioso de imitar, e está a empresa organizada para o aproveitar? Só um recurso que passe as quatro conta como vantagem competitiva \cite[p.~26]{PLAN}. Como o Relatório e Contas não compara a AstraZeneca com os seus rivais, as respostas abaixo são uma avaliação, não uma declaração da empresa.

\begin{xltabular}{\linewidth}{@{}N{2.4cm}Yccccp{2.5cm}@{}}
\caption{Análise VRIO: Valor, Raridade, Imitabilidade, Organização. O resultado segue a regra do material da unidade curricular \cite[p.~26]{PLAN}.}\label{tab:vrio}\\*
\toprule
\textbf{Recurso} & \textbf{Porquê} & \textbf{V} & \textbf{R} & \textbf{I} & \textbf{O} & \textbf{Resultado} \\*
\midrule
\endfirsthead
\toprule
\textbf{Recurso} & \textbf{Porquê} & \textbf{V} & \textbf{R} & \textbf{I} & \textbf{O} & \textbf{Resultado} \\*
\midrule
\endhead
\bottomrule
\endfoot
\textbf{Pipeline e I\&D} & 197 projetos; 20 medicamentos novos em fases avançadas ou em avaliação; 43 aprovações relevantes em 2025 nos EUA, UE, China e Japão \cite[pp.~28, 30]{AR} & S & S & S & S & \bgood{Vantagem} \\
\textbf{Liderança em oncologia} & \$25.6bn (+15\%) e cinco medicamentos acima de \$1bn; as afirmações de liderança são do próprio Grupo \cite[pp.~4, 12--14]{AR} & S & S & S & S & \bgood{Vantagem} \\
\textbf{Patentes dos produtos atuais} & A lei bloqueia as cópias enquanto duram; a Farxiga perde a exclusividade nos EUA em 2026 \cite[pp.~9, 33]{AR} & S & S & S & S & \bwarn{Até expirar} \\
\textbf{Excelência operacional} & Desempenho de fornecimento >99\% e 217 lançamentos dentro do prazo \cite[pp.~4, 27]{AR} & S & S & N & S & \bneut{Paridade} \newline{\scriptsize pode ser copiada} \\
\textbf{Presença global} & Vendas em mais de 125 países; os Mercados Emergentes são 26\% da receita \cite[pp.~2, 32--33]{AR} & S & N & N & S & \bneut{Paridade} \newline{\scriptsize os rivais também a têm} \\
\textbf{IA e dados} & Mais de 90\% da descoberta de moléculas pequenas usa IA; unidade Enterprise AI criada em dezembro de 2025; escalar a IA está listado como Risco Principal \cite[pp.~36, 48]{AR} & S & N & N & S & \bneut{Paridade} \newline{\scriptsize comum ao setor} \\
\textbf{Rede de parcerias} & Mais de 1,000 colaborações em todo o mundo \cite[p.~37]{AR} & S & N & N & S & \bneut{Paridade} \newline{\scriptsize os rivais também a têm} \\
\end{xltabular}

\begin{analysis}
A verdadeira vantagem da AstraZeneca é a sua \textbf{capacidade de inovar}: o pipeline, a I\&D e a sua posição em oncologia. As patentes também a protegem, mas apenas até expirarem. As operações, o alcance global, a IA e as parcerias são valiosos, mas os concorrentes também os têm ou os podem construir, pelo que não distinguem a empresa. Isto coincide com o que Porter mostrou: o sucesso depende de renovar o portefólio.
\end{analysis}

\subsection{Juntar tudo: a SWOT}
A SWOT reúne os pontos mais importantes das quatro ferramentas. Da PESTEL, mantêm-se os fatores com impacto alto ou médio; das outras ferramentas, os pontos que mais pesam ou que aparecem mais do que uma vez. Sob cada ponto, um rótulo cinzento mostra de que ferramenta veio, para que a origem da informação esteja sempre visível. As forças cujo rótulo está marcado como paridade passaram o teste do valor mas falharam o teste da raridade ou o da imitabilidade da VRIO: são necessárias para competir, mas não distinguem a empresa.

\begin{xltabular}{\linewidth}{@{}lp{3.6cm}Y@{}}
\caption{Matriz SWOT da AstraZeneca, com a ferramenta de origem de cada ponto.}\label{tab:swot}\\*
\toprule
\textbf{Código} & \textbf{Fator} & \textbf{Evidência e origem} \\*
\midrule
\endfirsthead
\toprule
\textbf{Código} & \textbf{Fator} & \textbf{Evidência e origem} \\*
\midrule
\endhead
\bottomrule
\endfoot
\multicolumn{3}{@{}l@{}}{\textbf{Forças --- dentro da empresa, hoje}} \\*
\midrule
S1 & \textbf{Pipeline e I\&D fortes} & \$14.2bn investidos (24\% da receita), 197 projetos, 20 medicamentos novos em fases avançadas \cite[pp.~28, 30, 125]{AR} \newline \textit{\small VRIO · Pipeline\quad Cadeia de valor · I\&D} \\
S2 & \textbf{Liderança em oncologia} & 44\% da receita, a crescer 15\% \cite[pp.~2, 4, 13]{AR} \newline \textit{\small VRIO · Oncologia} \\
S3 & \textbf{Alcance global, portefólio amplo} & Vendas em mais de 125 países; 16 medicamentos acima de \$1bn \cite[pp.~4, 32]{AR} \newline \textit{\small VRIO · Presença global: paridade\quad Cadeia de valor · Vendas} \\
S4 & \textbf{Operações fiáveis} & Fornecimento >99\%, 217 lançamentos dentro do prazo, zero achados críticos em 63 inspeções \cite[pp.~4, 27, 31]{AR} \newline \textit{\small VRIO · Operações: paridade\quad Cadeia de valor · Operações} \\
S5 & \textbf{Capacidade financeira e margem} & Margem bruta de 81.9\% e margem operacional de 23.4\%; fluxo de caixa operacional de \$14.6bn (+23\%), contra dívida líquida de \$23.4bn \cite[pp.~10, 61, 125]{AR} \newline \textit{\small Cadeia de valor · Finanças} \\
S6 & \textbf{Uso avançado de IA} & Mais de 90\% da descoberta de moléculas pequenas usa IA; unidade de IA dedicada \cite[pp.~5, 36]{AR} \newline \textit{\small VRIO · IA: paridade\quad Cadeia de valor · I\&D} \\
S7 & \textbf{Capacidade comprovada de reestruturar} & O programa PAAGR gerou benefícios anuais recorrentes antes de impostos, antes de reinvestimento, de \$1.9bn \cite[p.~59]{AR} \newline \textit{\small Cadeia de valor · Administração} \\
\midrule
\multicolumn{3}{@{}l@{}}{\textbf{Fraquezas --- dentro da empresa, hoje}} \\*
\midrule
W1 & \textbf{Demasiada receita em cinco produtos} & Farxiga, Tagrisso, Imfinzi, Ultomiris e Calquence somam \$30,045m: 51\% da Receita Total, ou 54\% das Vendas de Produtos em base comparável (calculado a partir de \cite[pp.~56, 125]{AR}) \newline \textit{\small Cadeia de valor · Vendas} \\
W2 & \textbf{Forte dependência dos EUA} & 43\% da receita, a caminho de cerca de 50\% até 2030 \cite[pp.~2, 4, 32]{AR} \newline \textit{\small Cadeia de valor · Vendas} \\
W3 & \textbf{Produtos mais antigos já a perder terreno} & Brilinta $-$38\% após os genéricos; Symbicort a erodir na UE e no Japão \cite[pp.~18, 20, 56]{AR} \newline \textit{\small Porter · Produtos substitutos\quad Cadeia de valor · Vendas} \\
W4 & \textbf{Crescimento lento ou negativo em algumas áreas} & Cardiovascular, renal e metabolismo +3\% (\$12.8bn, 22\% da receita); vacinas $-$13\% \cite[pp.~4, 17, 56]{AR} \newline \textit{\small Cadeia de valor · Vendas} \\
W5 & \textbf{Atrasada na meta climática de 2026} & $-$88.1\% de emissões face a uma meta de $-$98\%; o clima não é um Risco Principal no AR \cite[pp.~42, 47]{AR} \newline \textit{\small Cadeia de valor · Instalações} \\
W6 & \textbf{Contratempos no pipeline} & 20 projetos parados em 2025, 11 por segurança ou eficácia; a Fasenra falhou um ensaio de Fase III na DPOC \cite[pp.~20, 30]{AR} \newline \textit{\small Cadeia de valor · I\&D} \\
W7 & \textbf{Decisões operacionais concentradas no SET} & Todas as decisões operacionais significativas são tomadas ao nível do SET \cite[p.~147]{AR}, o que atrasa a resposta local de que a estratégia depende \newline \textit{\small Cadeia de valor · Administração} \\
\midrule
\multicolumn{3}{@{}l@{}}{\textbf{Oportunidades --- fora da empresa, futuro}} \\*
\midrule
O1 & \textbf{Procura crescente} & Envelhecimento e doença crónica \cite[pp.~6--7]{AR} \newline \textit{\small PESTEL · Social} \\
O2 & \textbf{Despesa em saúde a crescer} & +7.1\% por ano até 2029; vendas farmacêuticas +10\% em 2025 \cite[p.~6]{AR} \newline \textit{\small PESTEL · Económico} \\
O3 & \textbf{Mercados emergentes} & \$15.3bn (+12\%), +19\% fora da China \cite[pp.~4, 33]{AR} \newline \textit{\small PESTEL · Económico} \\
O4 & \textbf{IA e novas tecnologias} & \$4--7bn por ano para o setor; acesso através de mais de 1,000 colaborações \cite[pp.~7, 37]{AR} \newline \textit{\small PESTEL · Tecnológico} \\
O5 & \textbf{Grandes necessidades não satisfeitas} & Hipertensão, cancro e doenças raras, em que menos de 10\% têm tratamento \cite[pp.~13, 17, 23]{AR} \newline \textit{\small PESTEL · Social} \\
O6 & \textbf{Acordo com os EUA que premeia a produção local} & Isenção tarifária de três anos, a par de preços alinhados com os de outros países ricos \cite[pp.~4, 32--33]{AR} \newline \textit{\small PESTEL · Político} \\
O7 & \textbf{O setor é difícil de entrar} & Desenvolver um medicamento custa milhares de milhões e leva 5 a 15 anos, para além da barreira regulamentar, pelo que a ameaça de novos concorrentes é fraca e o valor fica com quem já está dentro \cite{MULLARD}, \cite[p.~9]{AR} \newline \textit{\small Porter · Novos concorrentes} \\
\midrule
\multicolumn{3}{@{}l@{}}{\textbf{Ameaças --- fora da empresa, futuro}} \\*
\midrule
T1 & \textbf{Pressão sobre os preços} & Orçamentos apertados, negociação IRA (Farxiga 2026, Calquence 2027), pagadores consolidados \cite[pp.~7, 32--33]{AR}, \cite{FEIN} \newline \textit{\small PESTEL · Económico, Legal\quad Porter · Clientes} \\
T2 & \textbf{Patentes a expirar} & A Farxiga (14.5\% da receita) perde a exclusividade nos EUA em 2026; o AR espera que isto seja gerível. Mais de \$90bn de vendas do setor em risco até 2029 \cite[pp.~33, 56]{AR}, \cite{IQVIA} \newline \textit{\small Porter · Produtos substitutos\quad VRIO · Patentes} \\
T3 & \textbf{Geopolítica e tarifas} & Rivalidade EUA--China e controlos ao comércio \cite[pp.~6, 49]{AR} \newline \textit{\small PESTEL · Político} \\
T4 & \textbf{Rivalidade intensa} & Grandes rivais de dimensão semelhante \cite{P50}, \cite[p.~3]{AR} \newline \textit{\small Porter · Rivalidade} \\
T5 & \textbf{Patentes mais fracas, novas regras da UE, atrasos nas aprovações} & \cite[pp.~48, 49]{AR}, \cite{EU} \newline \textit{\small PESTEL · Legal} \\
T6 & \textbf{Rutura no fornecimento e no clima} & Falta de medicamentos em níveis recorde, API externos concentrados na Índia e na China, fenómenos climáticos extremos \cite[pp.~7, 33, 44]{AR}, \cite{USP} \newline \textit{\small Porter · Fornecedores\quad PESTEL · Ambiental} \\
\end{xltabular}

\subsection{O que mais importa}
Os 27 pontos não pesam todos o mesmo. São ordenados pelo seu peso na receita, pela rapidez com que se fazem sentir, pelo facto de o Relatório e Contas os tratar ou não como prioridade ou Risco Principal, e pela possibilidade de serem medidos.

\medskip
\noindent\textbf{Dentro da empresa}
\begin{enumerate}[topsep=4pt,itemsep=2pt]
  \item 51\% da receita vem de cinco produtos, 14.5\% só da Farxiga; o outro lado do forte portefólio de oncologia \cite[pp.~2, 56]{AR}
  \item 9 dos 20 medicamentos novos prometidos para 2030 já foram entregues, sustentados por \$14.2bn de I\&D. O AR considera o sucesso da I\&D crucial \cite[pp.~10, 28, 52]{AR}
  \item \$14.6bn de fluxo de caixa operacional e um fornecimento >99\% financiam e executam o plano de \$50bn nos EUA, contra \$23.4bn de dívida líquida \cite[pp.~4, 10, 61]{AR}
\end{enumerate}

\medskip
\noindent\textbf{Fora da empresa}
\begin{enumerate}[topsep=4pt,itemsep=2pt]
  \item Em 2026 a Farxiga perde a exclusividade nos EUA e recebe um preço negociado no mesmo ano; a Calquence segue-se em 2027 \cite[pp.~33, 56]{AR}
  \item Passar de 43\% $\rightarrow$ 50\% da receita nos EUA até 2030 é a exposição da própria empresa (W2); o que está fora é a política que atua sobre ela --- preços negociados pelo Medicare e pagadores concentrados em três pharmacy benefit managers que tratam cerca de 80\% das receitas médicas dos EUA \cite[pp.~2, 4, 32]{AR}, \cite{FEIN}
  \item +7.1\% por ano na despesa em saúde e uma carga de doença crescente: o contrapeso às duas ameaças acima \cite[pp.~6, 7, 13, 17]{AR}
\end{enumerate}

\begin{analysis}
\textbf{A ligação entre os dois}\par
Tudo se resume a uma corrida: lançar medicamentos novos mais depressa do que os antigos perdem proteção. A Brilinta já perdeu 38\% das suas vendas e a Farxiga, o maior produto, está prestes a perder a exclusividade nos EUA \cite[pp.~33, 56]{AR}. O mercado continua a crescer em valor, mas cada produto rende durante menos tempo e a um preço mais baixo.
\end{analysis}

\subsection{Cruzamento dos fatores}
O material da unidade curricular sugere quatro perguntas para ligar o interior ao exterior \cite[pp.~27--28]{PLAN}. Cada célula responde a uma delas.

\begin{table}[H]\centering
\caption{Os quatro cruzamentos entre o interior e o exterior. As células não são dadas no Relatório e Contas.}
\label{tab:cruzamentos}
\begin{tabularx}{\linewidth}{@{}l>{\hsize=.5\hsize}Y>{\hsize=1.5\hsize}Y@{}}
\toprule
& \textbf{Cruzamento} & \textbf{O que significa para a AstraZeneca} \\
\midrule
\multirow{2}{*}{\rotatebox[origin=c]{90}{\textbf{\footnotesize Forças}}}
 & \textbf{Crescer} (SO) & Usar o pipeline e a oncologia (S1, S2) para responder à procura crescente, à despesa em saúde a subir e às necessidades não satisfeitas (O1, O2, O5), e a IA (S6) para acelerar a investigação (O4). As barreiras à entrada mantêm esses ganhos com as empresas que já estão dentro (O7). \\
\addlinespace
 & \textbf{Defender} (ST) & Usar o alcance global (S3) contra as tarifas e a geopolítica (T3), renovar o portefólio (S1) antes de as patentes expirarem (T2), e usar a liquidez e a capacidade comprovada de reestruturar (S5, S7) para financiar fábricas nos EUA e aquisições contra as tarifas (T3). \\
\midrule
\multirow{2}{*}{\rotatebox[origin=c]{90}{\textbf{\footnotesize Fraquezas}}}
 & \textbf{Corrigir} (WO) & Reduzir a dependência de cinco produtos (W1) entrando em áreas de grande necessidade como a hipertensão (O5), onde o baxdrostat tem dois resultados positivos de Fase III \cite[p.~17]{AR}. Os mercados emergentes (O3) acrescentam crescimento, mas com os mesmos produtos de topo \cite[p.~56]{AR}. \\
\addlinespace
 & \textbf{Evitar} (WT) & A concentração e os produtos em declínio (W1, W3) combinados com a pressão sobre os preços e as patentes a expirar (T1, T2) são a mistura mais perigosa, e o mesmo vale para a dependência dos EUA (W2) sob pressão de preços e de política (T1, T3). As decisões concentradas no SET (W7) atrasam a resposta local que a pressão sobre os preços (T1) exige. \\
\bottomrule
\end{tabularx}
\end{table}

\subsection{Estratégia proposta}
O material da unidade curricular pede que se ponderem várias alternativas antes de escolher uma \cite[pp.~3, 7, 29]{PLAN}. Não oferece uma lista de estratégias genéricas, e o Relatório e Contas não rejeita nenhuma opção, pelo que se comparam abaixo as três alavancas que o próprio relatório descreve.

\begin{table}[H]\centering
\caption{Alternativas consideradas antes de escolher a estratégia.}
\label{tab:alternativas}
\begin{tabularx}{\linewidth}{@{}lp{3.5cm}Y@{}}
\toprule
\textbf{Papel} & \textbf{Alavanca} & \textbf{Evidência e a que responde} \\
\midrule
Núcleo & Crescimento orgânico através da I\&D & Medicamentos novos e novas utilizações dos existentes; 12 primeiras aprovações de ciclo de vida em 2025 \cite[pp.~4, 52]{AR}. Responde a S1, S2 $\rightarrow$ O1, O5, O7, T2. \\
\addlinespace
Complemento & Aquisições, licenças e parcerias & Por exemplo a EsoBiotec, até \$1bn \cite[pp.~37, 48]{AR}. Responde a T2, W1, W3. \\
\addlinespace
Mantido como política & Remunerar os acionistas & Dividendo progressivo de \$3.20 por ação \cite[pp.~3, 10]{AR}. Importante, mas não corrige nenhuma fraqueza nem ameaça. \\
\bottomrule
\end{tabularx}
\end{table}
A estratégia proposta é o \textbf{crescimento através da inovação, com renovação mais rápida do portefólio, melhor equilíbrio geográfico e fornecimento seguro}. Baseia-se nos quatro cruzamentos, em cinco ações. O crescimento em si assenta em S1 e S2, os dois recursos que a VRIO confirma como vantagens que não expiram; S3, S4 e S6 são condições para executar a estratégia, não fontes de vantagem em si mesmas.

\begin{enumerate}[topsep=4pt,itemsep=3pt]
  \item \textbf{Focar a I\&D onde a empresa lidera ou quer liderar} Oncologia, cardiovascular, renal e metabolismo, respiratória e imunologia, e doenças raras, usando IA para reduzir o tempo de desenvolvimento. Medicamentos em fase avançada como o baxdrostat respondem também ao crescimento lento na área cardiovascular \cite[p.~17]{AR}. \textit{\small S1, S2, S6 $\rightarrow$ O1, O4, O5 · W4}
  \item \textbf{Substituir a receita que as patentes vão retirar} Com medicamentos novos e novas indicações, complementados com aquisições, licenças e parcerias \cite[pp.~37, 48]{AR}. \textit{\small S1, S5 $\rightarrow$ T2, W1, W3}
  \item \textbf{Crescer nos mercados emergentes e produzir localmente onde compensa} Os mercados emergentes trazem \$15.3bn e crescem 12\% \cite[pp.~4, 33]{AR}; fábricas locais nos EUA, na China e em Singapura \cite[pp.~4, 5, 32, 37]{AR}. Isto contraria a subida do peso dos EUA para cerca de 50\% até 2030 \cite[p.~4]{AR}. \textit{\small S3 $\rightarrow$ O3, O6, T3}
  \item \textbf{Mostrar aos pagadores quanto valem os medicamentos e deixar as regiões adaptar-se} Manter a política de preços em quatro patamares e os acordos baseados em resultados \cite[pp.~41, 48]{AR}, mas dar às equipas comerciais regionais margem para adaptar o acesso e o preço dentro de limites definidos centralmente. Hoje, todas as decisões operacionais significativas são tomadas ao nível do SET \cite[p.~147]{AR}; delegar esta aproxima a decisão do pagador que é preciso convencer, e o PAAGR mostra que o Grupo consegue reorganizar-se a essa escala \cite[p.~59]{AR}. \textit{\small S7 $\rightarrow$ T1 · W7}
  \item \textbf{Garantir o fornecimento e cumprir o plano climático} Manter o fornecimento acima de 99\% \cite[p.~4]{AR}, acrescentar capacidade \cite[p.~33]{AR} e fechar a distância para a meta de $-$98\% de emissões \cite[p.~42]{AR}. \textit{\small S4 $\rightarrow$ T6, W5}
\end{enumerate}

\begin{figure}[H]\centering
\begin{tikzpicture}[node distance=6mm and 6mm]
\node[fl] (c1) {Pipeline, I\&D e IA};
\node[fl, right=of c1] (c2) {Medicamentos novos, aquisições e licenças};
\node[fl, right=of c2] (c3) {Substituir a receita que perde proteção};
\node[fle, right=of c3] (c4) {Ambição 2030\\{\color{azmuted}$\geq$20 medicamentos novos, \$80bn}};
\node[fl, below=of c1] (d1) {Presença global e fábricas locais};
\node[fl, right=of d1] (d2) {Mercados emergentes, menos risco de tarifas};
\node[fl, right=of d2] (d3) {Valor para os pagadores, fornecimento seguro, plano climático};
\node[fle, right=of d3] (d4) {Crescimento sustentável};
\draw[se] (c1)--(c2); \draw[se] (c2)--(c3); \draw[se] (c3)--(c4);
\draw[se] (d1)--(d2); \draw[se] (d2)--(d3); \draw[se] (d3)--(d4);
\end{tikzpicture}
\caption{Como as cinco ações conduzem à Ambição 2030. As metas de 2030 são do Grupo \cite[p.~10]{AR}.}
\label{fig:acoes}
\end{figure}
Duas fraquezas ficam sem ação dedicada: a dependência dos EUA (W2) e os contratempos no pipeline (W6), que são aqui tratados como riscos a vigiar. A rivalidade e as regras das patentes (T4, T5) estão fora do controlo da empresa, e a estratégia responde-lhes indiretamente ao renovar o portefólio.

Como se compara isto com o que a AstraZeneca já diz? A sua estratégia declarada, \textit{Growth Through Innovation}, tem três prioridades \cite[pp.~2, 10, 11]{AR}, e o Relatório e Contas enumera mitigações semelhantes para os seus principais riscos \cite[p.~48]{AR}. As ações 1, 2 e 5 mantêm-se próximas dessa estratégia. As ações 3 e 4 afastam-se dela: a Ação 3 dá mais peso aos mercados emergentes do que a empresa dá, e a Ação 4 propõe delegar decisões de preço que hoje são tomadas centralmente. O que isto acrescenta é o caminho explícito da SWOT para os cruzamentos e destes para a ação, e a escolha entre alavancas.

\begin{xltabular}{\linewidth}{@{}N{2.5cm}Yp{2.2cm}Y@{}}
\caption{Indicadores para acompanhar a estratégia. As metas são declaradas pelo Grupo; a escolha dos indicadores não é. O Relatório e Contas não fixa nenhuma meta para a concentração da receita nem para as áreas de crescimento lento. A medida de fluxo de caixa do scorecard é o fluxo de caixa operacional menos o investimento em ativos fixos, mais as receitas de alienações de ativos intangíveis \cite[p.~91]{AR}, razão pela qual é mais baixa do que os \$14.6bn em S5.}\label{tab:indicadores}\\*
\toprule
\textbf{Objetivo} & \textbf{Indicador} & \textbf{Hoje} & \textbf{Meta declarada} \\*
\midrule
\endfirsthead
\toprule
\textbf{Objetivo} & \textbf{Indicador} & \textbf{Hoje} & \textbf{Meta declarada} \\*
\midrule
\endhead
\bottomrule
\endfoot
Renovar o portefólio & Medicamentos novos aprovados desde out. de 2022 & 9 & $\geq$ 20 até 2030 \cite[p.~10]{AR} \\
Crescer & Receita Total & \$58.7bn (2025) & \$80bn até 2030 \cite[p.~10]{AR} \\
Produtividade da I\&D & Eventos regulamentares por ano & 97 (2025) & Meta de 2025 de 54, com 69 contabilizados; 2026 não divulgado \cite[pp.~91, 102]{AR} \\
Produtividade da I\&D & Eventos de progressão do pipeline & 38 (2025) & Meta de 2025 de 29, com 36 contabilizados \cite[pp.~11, 91]{AR} \\
Capacidade financeira & Fluxo de caixa (medida do scorecard do Grupo) & \$11.9bn (2025) & Meta de 2025 de \$11.3bn \cite[p.~91]{AR} \\
Sustentabilidade & Emissões de Âmbito 1 e 2 face a 2015 & $-$88.1\% & $-$98\% até 2026 \cite[p.~42]{AR} \\
Sustentabilidade & Emissões de Âmbito 3 face a 2019 & --- & $-$50\% até 2030, $-$90\% até 2045; zero emissões líquidas até 2045 \cite[p.~42]{AR} \\
\end{xltabular}
